\section{Experimental Methodology}
\label{sec:methodology}

\subsection{Preprocessing}

Port numbers and protocol identifiers are categorical identifiers whose
magnitude is meaningless but whose value is informative, so alongside the raw
column we derive a log-scaled value, a well-known-port indicator, a
system-port indicator and an ephemeral-port indicator for each of the source
and destination ports. Three ratio features---bytes per packet, a burstiness
ratio of maximum to mean inter-arrival time, and packet-length range---are
added because they are row-wise deterministic and therefore safe to compute
before splitting. This yields $D = 27$ predictive columns for uni-flow and
$D = 40$ for bi-flow.

Every transformation with fitted parameters is fitted on the training split
alone: median imputation, winsorisation bounds taken from the 0.05th and
99.95th training percentiles, the quantile transform to an approximately normal
marginal, and the dropping of constant and near-zero-variance columns. The
preprocessing object refuses to transform before it has been fitted. The
columns \texttt{ip\_src}, \texttt{ip\_dst} and \texttt{is\_attack} are
rejected from the feature matrix by an assertion rather than by convention.

\subsection{Evaluation protocols}

Because the central question is what an evaluation can establish, we report
four protocols rather than one.

\emph{Protocol A, random stratified.} A 70/15/15 stratified split. This is the
protocol under which most of the literature, including the base paper, reports,
and we include it for comparability and as the object of study in its own
right.

\emph{Protocol B, duplicate-aware grouped.} Records whose exact predictive
feature vector coincides are grouped and assigned wholly to one split. Scan
traffic produces many identical records---11.4\% of uni-flow rows share a
feature vector with another row---so under Protocol A identical records
straddle the split and a model can be credited for retrieving a memorised
answer. Protocol B removes that route: it yields zero shared feature vectors
between training and test.

\emph{Protocol C, endpoint-grouped.} Whole endpoints are withheld, so the test
split contains devices the model has never seen. Which direction of holdout is
well posed follows from Section~\ref{sec:dataset}. Withholding \emph{sources}
cannot produce a test split containing both unseen sources and attack traffic,
because every attack class originates from one or two fixed addresses; such a
draw puts an attack class wholly in training or wholly in test. Withholding
\emph{victims} is well posed for the classes whose victims are distributed---
aggressive scan and UDP scan reach eighteen destinations each---and not for
MQTT brute-force (four victims) or Sparta SSH (two). Protocol C therefore
withholds victims, and we report the source-withholding variant separately as
Protocol CS to document the asymmetry rather than to conceal it.

A Protocol C or CS draw is accepted only if the test split contains at least
three classes, holds no more than 35\% of the records, and contains no class
absent from training; draws are redrawn under a recorded seed until one
qualifies, and the protocol reports itself \emph{not applicable} if none does.
Where a class is unavoidably absent from a test split, macro-averaged metrics
are averaged over the classes present and each run records which classes those
were, because averaging a structural zero over an absent class reports a
penalty the model never earned. Model selection uses the same convention as
evaluation, so early stopping optimises the quantity that is reported.

\emph{Protocol D, cross-representation.} Uni-flow and bi-flow are separate
feature spaces and are never merged; matched configurations are compared across
them.

\subsection{Baselines}

Twelve models are evaluated under identical splits.
\emph{EndpointLookup} is a majority-class lookup table on the ordered endpoint
pair, using no traffic features whatsoever; it exists to quantify how much of
the task is answerable from identity alone, and any graph model that merely
matches it has learned identity rather than behaviour. Random Forest,
XGBoost~\cite{chen2016xgboost}, SVM and $k$-nearest neighbours are the
classical learners; because a kernel SVM and a $k$-NN do not scale to a
350{,}000-row training split, both receive a stratified training cap whose size
is recorded with the run, and the SVM uses a Nystr\"om approximation of the RBF
kernel. GA-EBT is our reproduction of the base paper's detector. The MLP is
the graph-free deep baseline. GCN, GraphSAGE, E-GraphSAGE, GAT, EdgeGAT and
E-GAT-IDS are the graph family described in Section~\ref{sec:method}.

\subsection{GA-EBT reproduction fidelity}

We implement Algorithm~1 of the base paper literally: the chromosome encodes a
feature bit-mask together with the ensemble's number of estimators, maximum
depth and minimum samples per split; population 40, one hundred generations,
crossover 0.80, mutation 0.05, tournament size three, elitism of two, a
ten-feature minimum enforced after mutation, the published parameter bounds and
jitter magnitudes, and termination on generation count or ten generations
without fitness gain.

It remains an \emph{approximate} reproduction, and we state the deviations
rather than implying replication. We use a scikit-learn bagging ensemble of
decision trees where the original used WEKA, whose tree induction differs in
pruning and split criteria. We reduce the genetic search budget, because
$40 \times 100$ chromosomes evaluated by ten-fold cross-validation over a
350{,}000-row training split is not tractable on the hardware used here; the
budget was set empirically after a pilot search showed fitness reaching 0.99832
at generation zero and 0.99881 by generation five and then not improving for
ten further generations, which is unsurprising given that the fitness is
cross-validated F1 on a training split that is itself saturated. We run the
search once per (representation, task, protocol) and refit per seed. We apply
the original's min--max scaling inside this baseline so that it retains its own
published preprocessing. Three things cannot be reproduced at all: the
original's train/test partition, which is not published; its 33{,}079-instance
subset, whose sampling is not described; and the WEKA random seed. The complete
statement is carried in the code and in the run records.

\subsection{Metrics}

Accuracy is reported but is not the headline, because it saturates. The metrics
that carry the argument are macro-F1, the Matthews correlation
coefficient~\cite{chicco2020mcc}, PR-AUC~\cite{saito2015prcurve}, the
false-positive rate and minority-class recall. Balanced accuracy is reported
because under endpoint-grouped protocols it exposes a failure mode that macro-F1
compresses: a model that misses one of $k$ test classes entirely scores close to
$(k-1)/k$. Per-class precision, recall, F1 and false-positive rate are reported
for every run, together with training time, inference latency, throughput,
trainable parameter count and peak resident memory.

\subsection{Multi-seed protocol and statistics}

Every configuration is run with seeds 42, 123, 2024, 3407 and 2026. Results
are reported as mean $\pm$ standard deviation with 95\% bootstrap confidence
intervals over 10{,}000 resamples. Comparisons against the proposed model are
paired by seed, since every model sees an identical split for a given seed, and
we report the paired $t$-test, the Wilcoxon signed-rank test and a bootstrap
confidence interval on the mean difference, with Holm--Bonferroni correction
across the baseline family and Cohen's $d_z$ as effect size. We state
explicitly that at five paired observations the exact Wilcoxon test cannot
reach $p < 0.0625$, so significance claims rest on the confidence interval and
effect size as much as on the $p$-value; no test is run on a single number.

\subsection{Leakage controls}

The controls are enforced in code and recorded per run rather than asserted in
prose. Scaling, imputation, winsorisation and feature dropping are fitted on
training data only. Node attributes aggregate training flows only. Message
passing uses training edges only. Endpoint identifiers and the attack flag
cannot enter the feature matrix. Splits are verified for overlap and coverage.
The number of exact-duplicate feature vectors shared between training and test
is measured and reported rather than assumed absent. Early stopping and every
configuration decision use validation data; the test split is evaluated once,
after training concludes. An executable quality-control script re-checks these
properties over the stored artefacts and exits non-zero on failure.

\subsection{Hyperparameters and search}

We report results at the architectural defaults recorded in the configuration
files ($L = 2$ layers, hidden dimension 64, $H = 4$ heads, dropout 0.2,
AdamW at $10^{-3}$ with weight decay $10^{-5}$, cosine schedule, batch of 8192
target edges, early stopping on validation macro-F1 with patience 12 over at
most 60 epochs) rather than at the output of a large search. The reason is
stated rather than hidden: under Protocols A and B the validation metric is
saturated, so a search has no signal to follow; under Protocols C and CS the
differences we observe are architectural---between aggregation functions---and
an order of magnitude larger than anything the search space moves. A
successive-halving search over the specified grid is implemented and available,
and the grid, the selection metric and the budget are recorded.

\subsection{Reproducibility}

Every run writes a JSON record containing its metrics, per-class report,
confusion matrix, complete configuration, random seed, timings, parameter
count, peak memory, an environment fingerprint and the SHA-256 of every input
file. All tables and figures are generated from those records by script; no
number in this paper was transcribed by hand. Experiments ran on CPU only, and
the runtime figures should be read as such.
